% ============================================================================
%  1° Reporte de seguimiento (Agosto 2026) - Trabajo Terminal II
%
%  Compilación (desde esta carpeta):  pdflatex 1ReporteSeguimiento.tex
%
%  Las secciones citadas corresponden a Documentation/main.pdf
%  (reporte técnico de TT), que es el PDF que se adjunta como evidencia.
% ============================================================================
\documentclass[letterpaper,10pt,oneside]{article}
\usepackage[margin=2.54cm]{geometry}
\usepackage[spanish]{babel}
\newcommand\authormark[1]{\textsuperscript{#1}}
\usepackage[colorlinks=true,bookmarks=true,linkcolor=black,citecolor=black,urlcolor=blue]{hyperref}
\usepackage{times}
\usepackage[table]{xcolor}
\usepackage{booktabs}
\usepackage{graphicx}
\usepackage{titlesec}
\usepackage{tabularx}
\usepackage{array}
\usepackage{enumitem}
\usepackage{float}
\usepackage{xurl}
\urlstyle{same}

\titleformat{\section}
{\normalfont\fontsize{12pt}{14pt}\bfseries}
{\thesection}{1em}{}

% ----------------------------------------------------------------------------
%  Mes reportado
% ----------------------------------------------------------------------------
\def\numReporte{1}
\def\mes{Agosto}
\def\fechaCorte{31 de agosto de 2026}

% ----------------------------------------------------------------------------
%  Porcentajes de la Fase 3 (editar aquí).
%  "esp" = avance esperado según cronograma al corte, "real" = avance logrado.
% ----------------------------------------------------------------------------
\def\espEntrenamiento{33}   \def\realEntrenamiento{33}
\def\espExperimentos{33}    \def\realExperimentos{20}
\def\espAnalisis{33}        \def\realAnalisis{20}

% Color del encabezado de las tablas de evidencias
\definecolor{encabezado}{gray}{0.85}

% Porcentaje de cumplimiento: real respecto a lo esperado al corte
\newcommand\cumplimiento[2]{\fpeval{min(100, round(100*(#2)/(#1), 0))}\%}

% Porcentaje de la ficha: muestra el esperado solo si no se cumple
\newcommand\porcentaje[2]{#2\%\ifnum#2<#1\ (esperado: #1\%)\fi}

\begin{document}

	\begin{center}
		\fontsize{14}{16}\selectfont\bfseries
		\numReporte° Reporte de seguimiento (\mes{} 2026): \\
		Simulación de sistemas complejos mediante autómatas celulares neuronales
	\end{center}

	\begin{center}
		\fontsize{14}{16}\selectfont\bfseries
		\textit{Trabajo Terminal No. 2026-A114}
	\end{center}

	\begin{center}
		\fontsize{12}{16}\selectfont
		\textit{Alumno: \authormark{*}Angeles López Erick Jesse\authormark{1}} \\
		\textit{Director: Juárez Martínez Genaro\authormark{1}} \\
		\textit{\authormark{*}e-mail: eangelesl1700@alumno.ipn.mx}
	\end{center}

	\section{Información general}

	\begin{itemize}[leftmargin=*]
		\item \textbf{Periodo:} \mes{} 2026.
		\item \textbf{Evidencia:} reporte técnico de Trabajo Terminal. Capítulos 2 al 10.
		\item \textbf{Actividades:} Módulos de entrenamiento (cierre), Entrenamiento y ajuste, Experimentación y análisis, Análisis de resultados (global).
	\end{itemize}

	\section{Resumen de cumplimiento}

	\begin{table}[H]
		\centering
		\renewcommand{\arraystretch}{1.2}
		\begin{tabular}{l c c c c}
			\toprule
			\textbf{Actividad} & \textbf{Periodo} & \textbf{Esperado al corte} & \textbf{Avance real} & \textbf{Cumplimiento} \\
			\midrule
			\multicolumn{5}{l}{\textbf{Fase 1: Planificación y Revisión Bibliográfica}} \\
			Estado del arte                  & Ago/25 -- Oct/25 & 100\% & 100\% & 100\% \\
			Marco teórico y metodológico     & Sep/25 -- Oct/25 & 100\% & 100\% & 100\% \\
			\multicolumn{5}{l}{\textbf{Fase 2: Diseño e implementación del simulador}} \\
			Arquitectura del simulador       & Oct/25 -- Nov/25 & 100\% & 100\% & 100\% \\
			Módulos de simulación            & Nov/25 -- Dic/25 & 100\% & 100\% & 100\% \\
			Módulos entrenamiento            & Dic/25 -- Ago/26 & 100\% & 100\% & 100\% \\
			\multicolumn{5}{l}{\textbf{Fase 3: Entrenamiento y Validación}} \\
			Entrenamiento y ajuste           & Ago/26 -- Oct/26 & \espEntrenamiento\% & \realEntrenamiento\% & \cumplimiento{\espEntrenamiento}{\realEntrenamiento} \\
			Experimentación y análisis       & Ago/26 -- Nov/26 & \espExperimentos\%  & \realExperimentos\%  & \cumplimiento{\espExperimentos}{\realExperimentos} \\
			Análisis de resultados (global)  & Ago/26 -- Nov/26 & \espAnalisis\%      & \realAnalisis\%      & \cumplimiento{\espAnalisis}{\realAnalisis} \\
			\bottomrule
		\end{tabular}
	\end{table}

	\section{Evidencias por actividad}

	\subsection*{Fase 1: Planificación y Revisión Bibliográfica}

	\begin{table}[H]
		\centering
		\renewcommand{\arraystretch}{1.25}
		\begin{tabularx}{\linewidth}{|>{\raggedright\bfseries\arraybackslash}p{0.24\linewidth}|X|}
			\hline
			\rowcolor{encabezado}
			Nombre de la actividad & \textbf{Estado del arte} \\ \hline
			Periodo & Ago/25 -- Oct/25 \\ \hline
			Evidencia de trabajo & Documento escrito. \\ \hline
			Ubicación en el reporte & 5.1 (Estado del arte, aplicaciones y líneas actuales de investigación). \\ \hline
			Descripción del avance & Revisión de trabajos sobre autómatas celulares neuronales y sus aplicaciones. \\ \hline
			Porcentaje & 100\% \\ \hline
		\end{tabularx}
	\end{table}

	\begin{table}[H]
		\centering
		\renewcommand{\arraystretch}{1.25}
		\begin{tabularx}{\linewidth}{|>{\raggedright\bfseries\arraybackslash}p{0.24\linewidth}|X|}
			\hline
			\rowcolor{encabezado}
			Nombre de la actividad & \textbf{Marco teórico y metodológico} \\ \hline
			Periodo & Sep/25 -- Oct/25 \\ \hline
			Evidencia de trabajo & Documento escrito. \\ \hline
			Ubicación en el reporte & 2 (Sistemas, caos y simulación), 3 (Autómatas celulares), 4 (Inteligencia artificial y redes neuronales), 5 (Autómatas celulares neuronales) y 6 (Marco metodológico). \\ \hline
			Descripción del avance & Fundamentos teóricos del proyecto y diseño metodológico de la investigación. \\ \hline
			Porcentaje & 100\% \\ \hline
		\end{tabularx}
	\end{table}

	\subsection*{Fase 2: Diseño e implementación del simulador}

	\begin{table}[H]
		\centering
		\renewcommand{\arraystretch}{1.25}
		\begin{tabularx}{\linewidth}{|>{\raggedright\bfseries\arraybackslash}p{0.24\linewidth}|X|}
			\hline
			\rowcolor{encabezado}
			Nombre de la actividad & \textbf{Arquitectura del simulador} \\ \hline
			Periodo & Oct/25 -- Nov/25 \\ \hline
			Evidencia de trabajo & Documento escrito (análisis y diseño de software, diagramas UML). \\ \hline
			Ubicación en el reporte & 7 (Análisis de software), 8.1 (Arquitectura), 8.2 (Diagrama de clases), 8.3 (Diseño de datos) y 8.5 (Interfaces de usuarios). \\ \hline
			Descripción del avance & Requerimientos, casos de uso, diagramas UML, arquitectura MVC e interfaces gráficas. \\ \hline
			Porcentaje & 100\% \\ \hline
		\end{tabularx}
	\end{table}

	\begin{table}[H]
		\centering
		\renewcommand{\arraystretch}{1.25}
		\begin{tabularx}{\linewidth}{|>{\raggedright\bfseries\arraybackslash}p{0.24\linewidth}|X|}
			\hline
			\rowcolor{encabezado}
			Nombre de la actividad & \textbf{Módulos de simulación} \\ \hline
			Periodo & Nov/25 -- Dic/25 \\ \hline
			Evidencia de trabajo & Documento escrito, tablas de resultados y pruebas de operatividad. \\ \hline
			Ubicación en el reporte & 8.4 (Diseño de algoritmos), 9.1 (Entorno de desarrollo), 9.2 (Estructura del proyecto) y 9.4 (Validación y pruebas). Tablas 9.6 y 9.7. \\ \hline
			Descripción del avance & Simulador en C++/CUDA con interfaz en Qt, validado con reglas conocidas y pruebas de rendimiento. \\ \hline
			Porcentaje & 100\% \\ \hline
			Enlace de video & \url{https://drive.google.com/file/d/1tmUq2OIrWDl01cZxxLEUQqLCJ-hqWrWG/view?usp=drive_link} \\ \hline
		\end{tabularx}
	\end{table}

	\begin{table}[H]
		\centering
		\renewcommand{\arraystretch}{1.25}
		\begin{tabularx}{\linewidth}{|>{\raggedright\bfseries\arraybackslash}p{0.24\linewidth}|X|}
			\hline
			\rowcolor{encabezado}
			Nombre de la actividad & \textbf{Módulos entrenamiento} \\ \hline
			Periodo & Dic/25 -- Ago/26 \\ \hline
			Evidencia de trabajo & Documento escrito y pruebas de operatividad. \\ \hline
			Ubicación en el reporte & 7.8.4 (Entrenamiento), 8.4.5 (Entrenamiento del autómata celular neuronal), 9.3 (Entrenamiento del modelo) y 9.4.2 (Validación del autómata celular neuronal). \\ \hline
			Descripción del avance & Entrenamiento del NCA en PyTorch integrado a la interfaz gráfica. Concluido en el mes reportado. \\ \hline
			Porcentaje & 100\% \\ \hline
		\end{tabularx}
	\end{table}

	\subsection*{Fase 3: Entrenamiento y Validación}

	\begin{table}[H]
		\centering
		\renewcommand{\arraystretch}{1.25}
		\begin{tabularx}{\linewidth}{|>{\raggedright\bfseries\arraybackslash}p{0.24\linewidth}|X|}
			\hline
			\rowcolor{encabezado}
			Nombre de la actividad & \textbf{Entrenamiento y ajuste} \\ \hline
			Periodo & Ago/26 -- Oct/26 \\ \hline
			Evidencia de trabajo & Documento escrito y tablas de resultados. \\ \hline
			Ubicación en el reporte & 10.1 (Duplicación de reglas del juego de la vida). Tablas 10.1 y 10.2. \\ \hline
			Descripción del avance & Comparativa de funciones de activación y tamaños de red para la regla B3/S23. \\ \hline
			Porcentaje & \porcentaje{\espEntrenamiento}{\realEntrenamiento} \\ \hline
		\end{tabularx}
	\end{table}

	\begin{table}[H]
		\centering
		\renewcommand{\arraystretch}{1.25}
		\begin{tabularx}{\linewidth}{|>{\raggedright\bfseries\arraybackslash}p{0.24\linewidth}|X|}
			\hline
			\rowcolor{encabezado}
			Nombre de la actividad & \textbf{Experimentación y análisis} \\ \hline
			Periodo & Ago/26 -- Nov/26 \\ \hline
			Evidencia de trabajo & Documento escrito y gráficas de resultados. \\ \hline
			Ubicación en el reporte & 10.1 (Duplicación de reglas del juego de la vida). Figuras 10.1 y 10.2. \\ \hline
			Descripción del avance & Generación del conjunto de datos y análisis de su balance para la regla B3/S23. \\ \hline
			Porcentaje & \porcentaje{\espExperimentos}{\realExperimentos} \\ \hline
		\end{tabularx}
	\end{table}

	\begin{table}[H]
		\centering
		\renewcommand{\arraystretch}{1.25}
		\begin{tabularx}{\linewidth}{|>{\raggedright\bfseries\arraybackslash}p{0.24\linewidth}|X|}
			\hline
			\rowcolor{encabezado}
			Nombre de la actividad & \textbf{Análisis de resultados (global)} \\ \hline
			Periodo & Ago/26 -- Nov/26 \\ \hline
			Evidencia de trabajo & Documento escrito. \\ \hline
			Ubicación en el reporte & 10.1 (Duplicación de reglas del juego de la vida). \\ \hline
			Descripción del avance & Identificación del sesgo de muestra en patrones de vecindad poco frecuentes. \\ \hline
			Porcentaje & \porcentaje{\espAnalisis}{\realAnalisis} \\ \hline
		\end{tabularx}
	\end{table}

	\section{Actividades pendientes}

	\begin{table}[H]
		\centering
		\renewcommand{\arraystretch}{1.2}
		\begin{tabular}{l c l}
			\toprule
			\textbf{Fase} & \textbf{Periodo} & \textbf{Actividad} \\
			\midrule
			Fase 4: Documentación y Difusión & Oct/26 -- Dic/26 & Documentación técnica \\
			Fase 4: Documentación y Difusión & Oct/26 -- Dic/26 & Artículo académico \\
			\bottomrule
		\end{tabular}
	\end{table}

	\section{Cronograma}

	\begin{table}[H]
	\resizebox{\linewidth}{!}{%
		\begin{tabular}{p{5cm}*{10}{>{\centering\arraybackslash}p{1cm}}}
			\toprule
			\textbf{Actividad} & \textbf{Ago/25} & \textbf{Sep/25} & \textbf{Oct/25} & \textbf{Nov/25} & \textbf{Dic/25} & \cellcolor{black!70}\textcolor{white}{\textbf{Ago/26}} & \textbf{Sep/26} & \textbf{Oct/26} & \textbf{Nov/26} & \textbf{Dic/26} \\
			\midrule
			\multicolumn{11}{l}{\textbf{Fase 1: Planificación y Revisión Bibliográfica}} \\

			Estado del arte       & \cellcolor{gray!30} & \cellcolor{gray!30} & \cellcolor{gray!30} &  &  &  &  &  &  &   \\

			Marco teórico y metodológico &  & \cellcolor{gray!30} & \cellcolor{gray!30} &  &  &  &  &  &  &   \\

			\midrule

			\multicolumn{11}{l}{\textbf{Fase 2: Diseño e implementación del simulador}} \\
			Arquitectura del simulador &  &  & \cellcolor{gray!30} & \cellcolor{gray!30} &  &  &  &  &  &   \\

			Módulos de simulación  &  &  &  & \cellcolor{gray!30} & \cellcolor{gray!30} &  &  &  &  &   \\

			Módulos entrenamiento   &  &  &  &  & \cellcolor{gray!30} & \cellcolor{gray!30} &  &  &  &   \\

			\midrule

			\multicolumn{11}{l}{\textbf{Fase 3: Entrenamiento y Validación}} \\
			Entrenamiento y ajuste   &  &  &  &  &  & \cellcolor{gray!30} & \cellcolor{gray!30} & \cellcolor{gray!30} &   &   \\

			Experimentación y análisis  &  &  &  &  &  & \cellcolor{gray!30}  & \cellcolor{gray!30} & \cellcolor{gray!30} & \cellcolor{gray!30} &   \\

			Análisis de resultados (global) &  &  &  &  &  & \cellcolor{gray!30} & \cellcolor{gray!30} & \cellcolor{gray!30} & \cellcolor{gray!30} &   \\

			\midrule

			\multicolumn{11}{l}{\textbf{Fase 4: Documentación y Difusión}} \\
			Documentación técnica  &  &  &  &  &  &  &  & \cellcolor{gray!30}  & \cellcolor{gray!30} & \cellcolor{gray!30}  \\

			Artículo académico  &  &  &  &  &  &  &  & \cellcolor{gray!30}  & \cellcolor{gray!30}  & \cellcolor{gray!30}  \\

			\bottomrule
		\end{tabular}
	}
	\end{table}

	\section{Alumnos y directores}

	\vspace{1em}
	\begin{tabular}{@{} p{0.45\textwidth} p{0.45\textwidth} @{}}
		\textit{Erick Jesse Angeles López} & \textit{Dr. Genaro Juárez Martínez} \\[5em]
		Firma: \rule{6cm}{0.4pt} & Firma: \rule{6cm}{0.4pt} \\
	\end{tabular}

\end{document}
